\ifdefined\gkver\else\def\gkver{student}\fi
\documentclass[\gkver]{gaokaozhenti}
\begin{document}

\chapter{2003年全国理综旧课程}
%% paperType: 理综物理部分

\begin{enumerate}
\item
%% number: 17
%% paperName: 2003年全国理综旧课程第 17 题 6 分
%% typeId: 1
%% type: 单选题
%% chapter: 光学
%% point: 光的电磁说
%% method: 无
%% score: 6
%% degree: 850
%% duplicateId: 0
%% body:
一束单色光从空气射入玻璃中，则其\xzanswer{C}
\fourchoices[answer=C]
{频率不变，波长变长}
{频率变大，波长不变}
{频率不变，波长变短}
{频率变小，波长不变}
%% answer: C
%% memo:
\memoanswer{
光的频率由光源决定，光从空气射入玻璃时频率不变；光在玻璃中的传播速度小于在空气中的速度，
由 $v=\lambda f$ 可知，频率不变而波速减小时波长变短。
}

\item
%% number: 21
%% paperName: 2003年全国理综旧课程第 21 题 6 分
%% typeId: 1
%% type: 单选题
%% chapter: 电场
%% point: 电势能电势差电场力做功
%% method: 无
%% score: 6
%% degree: 450
%% duplicateId: 0
%% body:
图中虚线所示为静电场中的等势面 1、2、3、4，相邻的等势面之间的电势差相等，其中等势面 3 的电势为 0。一带正电的点电荷在静电力的作用下运动，经过 a、b 点时的动能分别为 $26\UeV$ 和 $5\UeV$。当这一点电荷运动到某一位置，其电势能变为 $-8\UeV$ 时，它的动能应为\xzanswer{C}
\onepicture[align=right]{figs/21.png}
\fourchoices[answer=C]
{$8\UeV$}
{$13\UeV$}
{$20\UeV$}
{$34\UeV$}
%% answer: C
%% memo:
\memoanswer{
点电荷从 a 到 b 动能减少 $21\UeV$，由能量守恒可知电势能增加 $21\UeV$。图中 a 在等势面 1 上、
b 在等势面 4 上，相邻等势面电势差相等，共跨 3 个间隔，故每相邻两等势面间的电势能差为
$7\UeV$。已知等势面 3 的电势为 0，则等势面 1 的电势为 $-14\UV$，a 点的电势能为 $-14\UeV$。
点电荷的总能量 $E=E_{k}+E_{p}=26\UeV+(-14\UeV)=12\UeV$。当电势能变为 $-8\UeV$ 时，
动能 $E_{k}'=E-E_{p}'=12\UeV-(-8\UeV)=20\UeV$。故选 C。
}

\item
%% number: 25
%% paperName: 2003年全国理综旧课程第 25 题 20 分
%% typeId: 6
%% type: 计算题
%% chapter: 交流电
%% point: 交流电
%% method: 无
%% score: 20
%% degree: 350
%% duplicateId: 0
%% body:
曾经流行过一种向自行车车头灯供电的小型交流发电机，图 1 为其结构示意图。图中 N、S 是一对固定的磁极，abcd 为固定在转轴上的矩形线框，转轴过 bc 中点、与 ab 边平行，它的一端有一半径 $r_{0}=1.0\Ucm$ 的摩擦小轮，小轮与自行车车轮的边缘相接触，如图 2 所示。当车轮转动时，因摩擦而带动小轮转动，从而使线框在磁极间转动。设线框由 $N=800$ 匝导线圈组成，每匝线圈的面积 $S=20\Ucmq$，磁极间的磁场可视做匀强磁场，磁感强度 $B=0.010\UT$，自行车车轮的半径 $R_{1}=35\Ucm$，小齿轮的半径 $R_{2}=4.0\Ucm$，大齿轮的半径 $R_{3}=10.0\Ucm$（见图 2）。现从静止开始使大齿轮加速转动，问大齿轮的角速度为多大才能使发电机输出电压的有效值 $U=3.2\UV$？（假定摩擦小轮与自行车轮之间无相对滑动）
\onepicture[w=12cm, align=right]{figs/25.png}
%% answer:
\jdanswer{
$3.2\Urads$。
}
%% memo:
\memoanswer{
设大齿轮、小齿轮、车轮、摩擦小轮的角速度分别为 $\omega_{3}$、$\omega_{2}$、$\omega_{1}$、$\omega_{0}$。
链条不打滑，大、小齿轮边缘线速度相等，有 $\omega_{3}R_{3}=\omega_{2}R_{2}$；小齿轮与自行车车轮共轴，
故 $\omega_{2}=\omega_{1}$；摩擦小轮与车轮边缘接触处不打滑，有 $\omega_{1}R_{1}=\omega_{0}r_{0}$。
由以上三式得 $\omega_{0}=\omega_{3}\frac{R_{3}R_{1}}{R_{2}r_{0}}$。

发电机线框产生的感应电动势的峰值 $E_{m}=NBS\omega_{0}$，输出电压的有效值
$U=\frac{E_{m}}{\sqrt{2}}=\frac{NBS\omega_{0}}{\sqrt{2}}$，故
\[
\omega_{3}=\frac{\sqrt{2}U r_{0}R_{2}}{NBS R_{3}R_{1}}
=\frac{\sqrt{2}\times3.2\times0.010\times0.040}{800\times0.010\times2.0\times10^{-3}\times0.10\times0.35}\Urads
\approx3.2\Urads.
\]
}

\end{enumerate}

\end{document}
